\section{Energy-exchanging boundary loads and front rigidity}
\label{sec:thermal_boundary}

A second test allows the thermal sector to exchange energy and connects
its response to the \emph{full} modular channel. It yields a positive
linear boundary model but a structural arithmetic exclusion. Complete
derivations, additional endpoint checks and the diagnostic record are
in \cite{BoundaryNote}.

\subsection{Native Kubo response and its physical scope}

Write $H_S=\sum_{\ell=2,3}\epsilon_\ell N_\ell$, with
$\epsilon_\ell=\log\ell$, and retain the native Gibbs state. The
isometry $v_\ell|m\rangle=|m+1\rangle$ obeys
$v_\ell^*v_\ell=I$ and $v_\ell v_\ell^*=I-P_\ell$.
Unlike the primitive gate, the bounded self-adjoint probe
$B_\ell=v_\ell+v_\ell^*$ changes the occupation energy. A real
boundary effort drives the interaction
\begin{equation}
 H_{\rm probe}(t)=H_S-U(t)\sum_\ell g_\ell B_\ell.
\end{equation}
For bounded real $U$, this is a bounded self-adjoint perturbation on
the original domain. All dynamical times use the common radiation time
conjugate to $p$; inverse temperature is a separate parameter.

The retarded response around the uncoupled Gibbs state is obtained
from the Kubo commutator \cite{Kubo}. Since
$[B_\ell(t),B_\ell(0)]=-2\ii\sin(\epsilon_\ell t)P_\ell$,
\begin{equation}
 \chi_{\ell,\beta}(t)=2(1-q_\ell)\sin(\epsilon_\ell t)1_{t>0},
 \qquad
 \chi_{\ell,\beta}(p)=\frac{2\epsilon_\ell(1-q_\ell)}
 {p^2+\epsilon_\ell^2}.
 \label{eq:native_kubo}
\end{equation}
The conjugate flow $I=\sum_\ell g_\ell\,d\langle B_\ell\rangle/dt$
has admittance
\begin{equation}
 Y_\beta(p)=\sum_{\ell=2,3}w_{\ell,\beta}
 \frac{p}{p^2+\epsilon_\ell^2},\qquad
 w_{\ell,\beta}=2g_\ell^2\epsilon_\ell(1-q_\ell)>0.
 \label{eq:native_admittance}
\end{equation}
It is positive real in $\HR$, because
$p/(p^2+\epsilon^2)=\tfrac12[(p-\ii\epsilon)^{-1}+
(p+\ii\epsilon)^{-1}]$. Here $\log\ell$ is a transition frequency,
not a derived flight delay.

For a common strength $\lambda\ge0$, this susceptibility has an exact
conservative linear realization:
\begin{align}
 \ddot x_\ell+\epsilon_\ell^2x_\ell
 &=\sqrt{\lambda w_{\ell,\beta}}\,U,
 & I&=\sum_\ell\sqrt{\lambda w_{\ell,\beta}}\,\dot x_\ell,\\
 E_{\rm load}&=\tfrac12\sum_\ell
 (|\dot x_\ell|^2+\epsilon_\ell^2|x_\ell|^2),
 &\dot E_{\rm load}&=\Ree(\bar U I).
 \label{eq:load_energy}
\end{align}
The initial coherent perturbation is zero. Equation~\eqref{eq:native_kubo}
is the exact susceptibility of the uncoupled probe;
\eqref{eq:load_energy} is an exact \emph{linear realization} of it.
Neither statement identifies the closed-loop model with exact
finite-coupling microscopic quantum scattering. That would require
an interacting equilibrium, radiative dressing and fluctuation channels.
Indeed
\begin{equation}
 \langle B_\ell(t)B_\ell(0)\rangle_\beta
 =e^{-\ii\epsilon_\ell t}+q_\ell e^{\ii\epsilon_\ell t}.
\end{equation}
The load energy balance concerns coherent perturbations, not the full
thermal output intensity of a noiseless quantum device.

\subsection{Connection to the completed modular response}

Use the established half-shift reference
$K(p)=\xi(p)/\xi(p+1)$ and $r_b(p)=e^{-bp}K(p)$ at an exterior port.
Its fixed-shift causality is an input from the known innerness result
and modular construction \cite{Suzuki,ModularNote}. The favorable
control $b=0$ removes the extra lead algebraically; it is not asserted
to be the selected regular geometric cross-section.

In the fixed measured-output orientation, let $f,g$ be external
amplitudes and $a,r_ba$ the core amplitudes. Impose common effort and
addition of flows:
\begin{equation}
 U=f-g=a-r_ba,\qquad f+g=a+r_ba+I,\qquad I=\lambda Y_\beta U.
\end{equation}
Then $f=a+I/2$, $g=r_ba+I/2$, and elimination yields
\begin{equation}
 S_{\beta,\lambda,b}
 =\frac{r_b+J_\beta(1-r_b)}{1+J_\beta(1-r_b)},\qquad
 J_\beta=\lambda Y_\beta/2.
 \label{eq:parallel_response}
\end{equation}
The entire completed modular transfer is retained. The core admittance
$(1+r_b)/(1-r_b)$ is positive real, so adding $\lambda Y_\beta$
and taking its Cayley transform proves that
\eqref{eq:parallel_response} is analytic and contractive in $\HR$.
For its numerator $N$ and denominator $D$, the flux identity is
\begin{equation}
 |D|^2-|N|^2=1-|r_b|^2+\lambda\Ree Y_\beta\,|1-r_b|^2\ge0.
\end{equation}
The ordinary finite-window coherent balance is
\begin{equation}
 \norm f_{(0,L)}^2-\norm g_{(0,L)}^2
 =E_{\rm load}(L)+\norm a_{(0,L)}^2-\norm{r_ba}_{(0,L)}^2\ge0.
 \label{eq:parallel_energy}
\end{equation}
Here $r_ba$ denotes causal convolution. The core deficit is the later
output energy for its input stopped at $L$. This uses the established
half-shift isometry, not positivity of the unknown variable-shift target.

\subsection{A structural exclusion for finite spectral mass}

\begin{proposition}[Fixed-core parallel-load front rigidity]
\label{prop:front_rigidity}
Suppose $K(p)\sim c\,p^{-1/2}$, $c>0$, on the positive real axis,
and a load has the representation
\begin{equation}
 Y(p)=\int_{(0,\infty)}\frac{p}{p^2+\Omega^2}\dd\nu(\Omega),
 \qquad \nu\ge0,\qquad C:=\nu((0,\infty))<\infty.
 \label{eq:finite_spectral_mass}
\end{equation}
For each fixed finite $\lambda\ge0$, the connection
\eqref{eq:parallel_response} at $b=0$ satisfies
\begin{equation}
 S(p)=K(p)+\frac{\lambda C}{2p}+o(p^{-1}),\qquad
 S(p)\sim c\,p^{-1/2}.
 \label{eq:front_rigidity}
\end{equation}
For the modular core, it cannot equal $K_{\beta/2}$ at any
$0<\beta<1$.
\end{proposition}
\begin{proof}
Dominated convergence gives
$pY(p)=\int p^2/(p^2+\Omega^2)\dd\nu(\Omega)\to C$.
Subtracting the core from \eqref{eq:parallel_response} gives
\begin{equation}
 S-K=\frac{(\lambda Y/2)(1-K)^2}{1+(\lambda Y/2)(1-K)}.
\end{equation}
This proves \eqref{eq:front_rigidity}. The required leading behavior is
$K_{\beta/2}(p)\sim(2\pi/p)^{\beta/2}$, which has a different power
for $0<\beta<1$.
\end{proof}

The hypothesis includes countably many modes when their total response
spectral mass is finite. For a self-adjoint thermal probe $B$, the
transition measure has positive weights
$2(E_n-E_m)(\rho_m-\rho_n)|B_{mn}|^2$ at $E_n-E_m>0$.
When the sum converges its total mass is
\begin{equation}
 C_\beta=\operatorname{Tr}\rho_\beta[B,[H_S,B]].
\end{equation}
Finite sums of products of the native isometries and their adjoints
have bounded energy jumps and the needed bounded commutators. In
particular, \eqref{eq:native_admittance} has
$C_\beta=\sum_\ell w_{\ell,\beta}<\infty$, without truncating the
occupation space. Boundedness of an arbitrary $B$ alone does not
imply this energy regularity for an unbounded Hamiltonian.

At a regular exterior port $b>0$, a nonzero load instead gives
$S_{\beta,\lambda,b}(p)\sim\lambda C_\beta/(2p)$: it introduces a
prompt response before the reference delay of
$e^{-bp}K_{\beta/2}$. Thus the regular port does not repair the failure
already visible in the favorable $b=0$ comparison.

\subsection{The calibrated tangent and endpoints also fail}

A fixed active added load changes the half-shift reference itself.
Grant the favorable additional switch $\lambda(\beta)=1-\beta$ so
that $S_{1,0,0}=K$. This switch is not supplied by the KMS law.
Its microscopic amplitude is proportional to $\sqrt{1-\beta}$;
a differentiable amplitude vanishing at one would give zero first
response tangent. The one-sided tangent under the granted switch is
\begin{align}
 T_{\rm native}(p)
 &=\left.\partial_\beta\log S_{\beta,1-\beta,0}(p)\right|_1
 =-\frac{Y_1(p)(1-K(p))^2}{2K(p)}
 \sim-\frac{C_1}{2\sqrt{2\pi}}p^{-1/2},\\
 T_{\rm arith}(p)
 &=\left.\partial_\beta\log K_{\beta/2}(p)\right|_1
 =-\tfrac12 a_{1/2}(p)
 =-\tfrac12\log\frac{p}{2\pi}+O(p^{-1}).
 \label{eq:boundary_tangent_failure}
\end{align}
For $g_2=g_3=1$, $C_1=\log2+(4/3)\log3$. No constant calibration
can turn the decaying native tangent into the required logarithm.
In time, the native transfer tangent has a finite step at the front,
whereas the arithmetic transfer tangent has a
$t^{-1/2}\log t$ singularity \cite[Section~7]{BoundaryNote}.

For the explicit finite-transition load, $Y_\beta(p)=O(p)$ near zero,
so $S'_{\beta,\lambda,0}(0)=K'(0)$, while the target slope varies
with $\beta$. For bounded $\lambda(\beta)$, the native weights vanish
as $\beta\downarrow0$, giving $S\to K$ pointwise rather than the
identity. These supplementary failures are not needed for the
front-rigidity proof.

The resulting decision is to park the regular finite-place passive
parallel-load route. The theorem concerns the unchanged core and
\eqref{eq:finite_spectral_mass}; it excludes neither distributed
interactions that change the core nor every passive boundary admittance.
A singular continuum proposal would need its own well-defined domain
and positive energy. The coefficient identity in
\eqref{eq:main_dictionary} remains intact.
